\documentclass[12pt]{article}

\begin{document}
My project, with partner Eric Hu, seeks to characterize the emissivity of a material as a function of its finish, and we seek to obtain this value through the numerical approximation of the characteristic cooling curve of the said material. Here, our independent variable is the finish, and the dependent variable is the emissivity. To do this, we utilize three uniform aluminum pipes (of roughly the same height, diameter, and thickness), of rough, smooth, and lacquered finishes each, heat all three pipes for the same period of time from the same lab temperature, then measure the temperature of each of the pipes at consistent times. This summarizes our data collection, and from theoretical expectations, all distinct cooling curves have a characteristic emisivity, so we intend to obtain the emissivity from this data by fitting a heat decay model to the data in the form of the numerically solved conduction-less differential temperature model. We utilize a conduction-less model as we assume that conduction is minimal due to the exposure of the pipes to open air and the high conductivity of aluminum (so there is a near-zero temperature gradient across the alumium pipe itself). Over the past week's lab, we were able to obtain a set of measurements necessary for fitting the pipes alongside a set of test data for testing the differential equation solver we would use for future experiments, but we were unable to actually complete the differential equation solver in time to obtain an emissivity.

In the next lab, I hope that we can setup an automatic temperature measurement system to complete our experiments in parallel and to finish the temperature modelling system to obtain emissivity values from characteristic cooling curves. Firstly, to maximize the efficiency in our time (as only three hours in lab is not much), a system to automatically take temperature data in small time increments with automatically-calculated uncertainties would free up hours of annoying manual data collection and management work that can be spent performing deeper analysis and performing calculations. Next, finishing the differential equation solver will allow us to begin obtaining emissivity values and comparing them with literature (which can allow us to ask further, deeper questions about the accuracy of our findings).

In terms of improvements, the biggest improvement to implement in our next experimental cycle is to utilize thermocouple thermometers for temperature measurement over visual analysis to reduce uncertainty in measurement (due to a phone measurement being inaccurate due to human reading and the coriolis effect from the camera perspective on the thermometer) and parallelizing our data collection through the use of multiple-output thermocouples as opposed to having to take a singular dataset at a time. 
\end{document}
